\chapter{Introduction}
\label{chap:introduction}
[Introduce the topic, explain the research problem and state the contribution
of the thesis. All prose and references in this example are placeholders.]

\section{Research aims and questions}
[Explain the overall aim and formulate the research questions.]
\begin{enumerate}
  \item {}[First research question.]
  \item {}[Second research question.]
\end{enumerate}

\section{Comments}

During drafting, you can use comments to flag unclear ideas, leave
reminders, and exchange feedback with your supervisors. Named comment
commands make it easy to see who added each note.

\kristiina{This is a fuzzy idea.}
\supervisor{Please elaborate.}

The comment commands are defined in \texttt{commands.tex} using the
\texttt{todonotes} package. You can add more commands or change the
names and colours to suit your collaborators.

Before submitting the thesis, remove the comments or hide them by
replacing the package-loading line in \texttt{commands.tex} with
\verb|\usepackage[disable]{todonotes}|.


\section{Structure of the thesis}
[Explain the relationship between the chapters or included publications.]
Chapter~\ref{chap:methods} demonstrates a table, figure and linguistic example.

\subsection*{Citation examples}
The following citations refer to explicitly fictional demonstration records.
A narrative citation uses \citet{examplebook} or \citet[22]{examplebook}; a parenthetical citation uses
\citep{examplearticle} or \citep[10]{examplearticle}. A page reference can be written as
\citep[12--14]{examplebook} or with multiple references \citetext{\citealp[30--44]{examplebook}; \citealp{examplearticle}}.

\subsection*{Footnotes and quotations}
A short quotation can use \enquote{language-aware quotation marks}.
A footnote is inserted with the ordinary LaTeX command.\footnote{Example
footnote. Replace with a relevant note or remove it.}
